%%%% rapport.tex --- INF8225 Projet 2026
\typeout{INF8225 Projet -- Pipelines hybrides pour la segmentation de tumeurs pancreatiques}

\documentclass{article}
\pdfpagewidth=8.5in
\pdfpageheight=11in
\usepackage{ijcai19}

\usepackage{times}
\usepackage{soul}
\usepackage{url}
\usepackage[hidelinks]{hyperref}
\usepackage[utf8]{inputenc}
\usepackage[small]{caption}
\usepackage{graphicx}
\usepackage{amsmath}
\usepackage{amssymb}
\usepackage{booktabs}
\urlstyle{same}

\title{De la d\'etection de polypes \`a la segmentation de tumeurs pancr\'eatiques :\\
limites des d\'etecteurs \emph{open-vocabulary} et r\^ole d'un classifieur appris en aval}

\author{
Morad Bel-Melih$^{1}$\footnote{Auteur de contact}\and
Paul B.$^{1}$\\
\affiliations
$^{1}$Polytechnique Montr\'eal, Cours INF8225 -- Projet 2026\\
\emails
morad.bel-melih@polymtl.ca,
paul.b@polymtl.ca
}

\begin{document}

\maketitle

\begin{abstract}
Nous \'etudions une cha\^ine de segmentation faiblement supervis\'ee qui combine
un d\'etecteur \emph{open-vocabulary} (\emph{Grounding DINO}) et un segmenteur
\emph{promptable} (\emph{MedSAM}). Cette cha\^ine atteint un DICE de 0{,}95
sur \emph{Kvasir-SEG} en mode oracle, ce qui valide le principe d'une
segmentation pilot\'ee par bo\^ites textuelles. Lorsque nous appliquons la
m\^eme architecture aux tumeurs pancr\'eatiques du \emph{Medical Segmentation
Decathlon}, la sp\'ecificit\'e s'effondre \`a 0{,}17 \`a cause des hallucinations
massives du d\'etecteur sur les coupes saines. Nous documentons ensuite quatre
strat\'egies successives pour restaurer cette sp\'ecificit\'e, motiv\'ees par
nos observations exp\'erimentales : un agent multimodal de v\'erification
(\emph{Gemini}), une porte anatomique fond\'ee sur un U-Net du pancr\'eas,
une cascade spatiale pancr\'eas$\rightarrow$tumeur, et enfin un classifieur
ResNet-18 d'ensemble entra\^in\'e sur les n\'egatifs durs produits par
\emph{Grounding DINO}. Cette derni\`ere strat\'egie porte la sp\'ecificit\'e
\`a 0{,}74 sans effondrer la sensibilit\'e, et illustre que le signal de
texture local capt\'e par un classifieur appris est compl\'ementaire au signal
\emph{open-vocabulary} du d\'etecteur.
\end{abstract}

\section{Introduction}

L'apprentissage profond a profond\'ement modifi\'e la segmentation d'images
m\'edicales. Les architectures sp\'ecialis\'ees comme \emph{U-Net}~\cite{ronneberger2015unet}
puis \emph{nnU-Net}~\cite{isensee2021nnunet} dominent les
\emph{benchmarks} lorsqu'un volume suffisant de masques pixel pr\'ecis est
disponible. Cette exigence reste un frein op\'erationnel, et plusieurs travaux
r\'ecents proposent de la contourner en composant deux mod\`eles de fondation.
\emph{Segment~Anything}~\cite{kirillov2023sam} et son adaptation m\'edicale
\emph{MedSAM}~\cite{ma2024medsam} produisent un masque \`a partir d'un
\emph{prompt} g\'eom\'etrique, tandis que les d\'etecteurs
\emph{open-vocabulary} comme \emph{Grounding~DINO}~\cite{liu2024groundingdino}
g\'en\`erent ce \emph{prompt} \`a partir d'une simple cha\^ine de caract\`eres.
La promesse est s\'eduisante~: aucune annotation pixel n'est requise au moment
de l'inf\'erence, seulement le nom de la classe d'int\'er\^et.

Notre projet \'eprouve cette cha\^ine sur deux probl\`emes de difficult\'e
contrast\'ee. Nous v\'erifions d'abord son comportement sur \emph{Kvasir-SEG},
un \emph{benchmark} de polypes coliques o\`u les objets pr\'esentent un fort
contraste visuel et des contours nets. Nous l'appliquons ensuite aux tumeurs
pancr\'eatiques 2D extraites du \emph{Medical Segmentation
Decathlon}~\cite{antonelli2022msd}, un cas notoirement difficile o\`u la
tumeur partage la texture du parenchyme environnant. Le passage du premier au
second domaine fait appara\^itre un mode d'\'echec sp\'ecifique, l'hallucination
massive de \emph{DINO} sur les coupes saines, qui m\`ene \`a une sp\'ecificit\'e
proche du hasard. La majeure partie de notre d\'emarche consiste alors \`a
identifier la source de ce mode d'\'echec et \`a tester quatre m\'ecanismes de
correction de complexit\'e croissante. Le pipeline final, qui ajoute un
classifieur ResNet-18 d'ensemble~\cite{he2016resnet} en garde-fou de second
\'etage, restaure une sp\'ecificit\'e exploitable. Le code et les donn\'ees
pr\'epar\'ees sont publi\'es \`a~\url{https://github.com/moradBMH/INF8225_Projet}.

\section{Travaux Ant\'erieurs}

La segmentation pancr\'eatique sp\'ecialis\'ee a \'et\'e \'etudi\'ee
intensivement~; Zhou \emph{et~al.}~\shortcite{zhou2017fixedpoint} introduisent
un \emph{fixed-point model} qui raffine it\'erativement la pr\'ediction de
l'organe, et Zhu \emph{et~al.}~\shortcite{zhu2019multi} ajoutent une cascade
multi-\'echelle pour le d\'epistage des ad\'enocarcinomes ductaux. Ces approches
exigent toutefois un volume cons\'equent de masques pixel-pr\'ecis 3D,
contrainte que \emph{nnU-Net}~\cite{isensee2021nnunet} a optimis\'ee mais sans
la lever.

L'arriv\'ee de \emph{SAM}~\cite{kirillov2023sam} a chang\'e la donne en
proposant un segmenteur g\'en\'eraliste pilot\'e par un \emph{prompt}
g\'eom\'etrique. \emph{MedSAM}~\cite{ma2024medsam} sp\'ecialise ce mod\`ele sur
1{,}5~M de paires image/masque m\'edicales et obtient des performances
comp\'etitives sur de nombreuses modalit\'es. La question reste de savoir
comment fournir cette bo\^ite automatiquement, et \emph{Grounding
DINO}~\cite{liu2024groundingdino} y r\'epond en alignant un encodeur visuel
DINO~\cite{caron2021dino} avec un encodeur textuel BERT~\cite{devlin2019bert}.
Plusieurs publications r\'ecentes assemblent ces deux composants pour des
\emph{benchmarks} m\'edicaux, mais peu \'evaluent honn\^etement leur
sp\'ecificit\'e sur des images sans la classe cible.

L'utilisation d'agents multimodaux pour raffiner les pr\'edictions visuelles
constitue une famille \'emergente de techniques inspir\'ees du \emph{Chain-of-Thought}
\cite{wei2022chainofthought}. Plusieurs travaux interrogent un LLM
multimodal comme \emph{Gemini}~\cite{team2024gemini} \`a propos d'une r\'egion
candidate. Notre tentative dans ce sens, document\'ee en section~\ref{sec:agent},
\'eclaire les limites pratiques de cette approche dans un contexte radiologique.

\section{Approche Th\'eorique}

\subsection{Principe d'une segmentation \emph{prompt\'ee}}

\emph{MedSAM} utilise l'architecture \emph{ViT-B} de \emph{SAM} et produit un
masque $\hat{\boldsymbol{M}}$ en moins de 50~ms par image, conditionnellement
\`a une bo\^ite $\boldsymbol{b}\in\mathbb{R}^4$ fournie en entr\'ee. La
qualit\'e du masque d\'epend essentiellement de la qualit\'e de la bo\^ite~;
toute erreur en amont se r\'epercute directement en aval, ce qui fait du
d\'etecteur le point critique du pipeline.

\emph{Grounding DINO} aligne un encodeur Swin-T avec un encodeur BERT par
\emph{cross-attention} bidirectionnelle. Pour une requ\^ete textuelle
$\boldsymbol{q}$ et une position spatiale $\boldsymbol{p}$, il pr\'edit un
score $s(\boldsymbol{p},\boldsymbol{q})\in[0,1]$ et une bo\^ite candidate.
La perte de classification est une \emph{focal loss}~\cite{lin2017focal} avec
$\gamma=2{,}0$ et $\alpha=0{,}25$, qui att\'enue le poids des nombreux ancres
faciles. La proc\'edure d'\'echantillonnage \emph{RandomSamplingNegPos}
introduit des cha\^ines \emph{distractor} pendant l'entra\^inement, ce qui
permet en principe au mod\`ele d'apprendre \`a ne rien activer en l'absence
de la classe cible.

\subsection{Apprentissage d'un \emph{prior} anatomique}

Pour fournir une r\'egion d'int\'er\^et anatomique, nous entra\^inons un
\emph{U-Net}~\cite{ronneberger2015unet} \`a quatre niveaux ($\sim$1{,}9~M
param\`etres) sur le label \texttt{1} (pancr\'eas) du MSD. Comme cette
porte doit avoir une sensibilit\'e \'elev\'ee pour ne pas masquer une vraie
tumeur, nous adoptons une perte de Tversky~\cite{salehi2017tversky} qui
p\'enalise davantage les faux n\'egatifs que les faux positifs,
\begin{equation}
\mathcal{L}_{\text{Tv}}=1-\frac{TP}{TP+\alpha\,FP+\beta\,FN},
\end{equation}
avec $\alpha=0{,}3$ et $\beta=0{,}7$. Le masque pr\'edit est dilat\'e de 35~px
pour absorber les erreurs de contour.

\subsection{Classifieur appris en second \'etage}

Lorsque les filtres spatiaux et s\'emantiques se r\'ev\`elent insuffisants
(section~\ref{sec:exp}), nous ajoutons un classifieur appris qui juge
chaque candidat localement. Concr\`etement, un \emph{crop} 224$\times$224
autour de chaque bo\^ite est pass\'e \`a cinq ResNet-18~\cite{he2016resnet}
initialis\'es avec des graines diff\'erentes. Le score \emph{image-level}
agr\`ege les probabilit\'es individuelles avec une p\'enalit\'e sur la
variance,
\begin{equation}
\hat{p}_{\text{ens}}=\overline{\hat{p}}-0{,}25\,\sigma(\hat{p}),
\end{equation}
ce qui filtre les candidats sur lesquels les mod\`eles sont en d\'esaccord.

\section{Donn\'ees et Protocole}

\emph{Kvasir-SEG} contient 1000 images d'endoscopie color\'ee accompagn\'ees
de masques de polypes et de bo\^ites \emph{ground-truth}. Nous l'utilisons
comme banc d'essai initial pour la cha\^ine \emph{detection$+$segmentation},
avec un \emph{split} 880/20/100 \`a la mani\`ere des publications standard.

\emph{MSD~Pancreas} contient 281 volumes 3D de scanners CT. Nous extrayons
des coupes axiales 2D et obtenons 800 images d'entra\^inement (400 contenant
une tumeur, 400 sans), 100 de validation et 100 de test, chacune \'equilibr\'ee
50/50 entre coupes positives et n\'egatives. Les masques sont multiclasse,
avec \texttt{0} pour le fond, \texttt{1} pour le pancr\'eas et \texttt{2} pour
la tumeur. Toutes les m\'etriques de segmentation sont calcul\'ees sur la
classe tumeur, et les m\'etriques de classification (sensibilit\'e,
sp\'ecificit\'e, F1) sur l'image enti\`ere.

\section{Exp\'eriences}\label{sec:exp}

Nous avons men\'e nos exp\'eriences en suivant une logique de complexit\'e
croissante~: chaque it\'eration r\'epond \`a une faiblesse mesur\'ee \`a
l'\'etape pr\'ec\'edente. Cette section retrace cette progression.

\subsection{Validation initiale sur Kvasir-SEG}

Avant d'aborder le probl\`eme pancr\'eatique, nous v\'erifions que la cha\^ine
\emph{detection$+$MedSAM} fonctionne sur un cas favorable. Sur \emph{Kvasir-SEG},
les polypes pr\'esentent un fort contraste avec la muqueuse environnante et
des contours bien d\'efinis. Nous \'evaluons d'abord MedSAM en mode oracle, en
lui fournissant directement les bo\^ites \emph{ground-truth} fournies avec
le dataset. La sortie atteint un DICE moyen de \textbf{0{,}954} sur les 1000
images, ce qui confirme que MedSAM est un segmenteur fiable lorsque la
bo\^ite est correcte. Le \emph{fine-tuning} de \emph{Grounding~DINO} sur la
classe \texttt{polyp} (configuration \texttt{polyp\_config.py}) permet ensuite
de fermer la boucle automatiquement, avec un mAP\(_{\text{COCO}}\) sup\'erieur
\`a 0{,}55 sur la validation. Ces r\'esultats valident le principe et nous
encouragent \`a porter la m\^eme architecture sur un probl\`eme plus
ambitieux.

\subsection{Premi\`ere application au pancr\'eas}

Le passage \`a \emph{MSD~Pancreas} introduit deux difficult\'es. D'abord, la
classe \texttt{tumor} n'appartient pas au vocabulaire pertinent du
\emph{checkpoint} pr\'e-entra\^in\'e \texttt{swin-T objects365$+$goldg}~;
appliqu\'e en \emph{zero-shot}, le mod\`ele produit des bo\^ites quasiment
al\'eatoires (mAP $<$ 0{,}05). Nous adaptons donc \emph{Grounding~DINO}
\`a deux classes (\texttt{pancreas}, \texttt{tumor}) en d\'egelant les quatre
\emph{stages} du backbone, en abaissant le taux d'apprentissage \`a $10^{-5}$,
en activant trois cha\^ines \emph{distractor} \'echantillonn\'ees, et en
utilisant le \emph{prompt} compos\'e ``\texttt{pancreas . tumor .}''.
L'entra\^inement converge en 25~\'epoques sur GPU~T4 et le meilleur
\emph{checkpoint} atteint un mAP\(_{\text{COCO}}\) d'environ 0{,}30 sur la
validation, ce qui constitue notre meilleur d\'etecteur.

Une fois le d\'etecteur sp\'ecialis\'e, nous reprenons la cha\^ine canonique
en conservant la bo\^ite tumeur de plus haut score et en l'envoyant \`a
\emph{MedSAM}. Le tableau~\ref{tab:abl} reporte les m\'etriques sur le
\emph{split} test. La sensibilit\'e est quasi-parfaite (0{,}98), mais la
sp\'ecificit\'e tombe \`a 0{,}17 parce que \emph{DINO} produit
syst\'ematiquement une bo\^ite m\^eme sur les coupes sans tumeur. Comme
\emph{MedSAM} segmente toujours quelque chose dans cette bo\^ite, la
sp\'ecificit\'e refl\`ete directement le comportement du d\'etecteur. Le
DICE moyen sur les images positives reste mod\'er\'e (0{,}47), une partie
des bo\^ites \'etant centr\'ee sur du parenchyme sain plut\^ot que sur la
tumeur.

\begin{table}[t]
\centering
\small
\setlength{\tabcolsep}{4pt}
\begin{tabular}{lrrrr}
\toprule
M\'ethode & Sens. & Sp\'ec. & DICE & F1 \\
\midrule
Oracle MedSAM (bo\^ite GT)              & 1{,}00 & --     & 0{,}88 & --     \\
DINO+MedSAM \emph{one-shot}             & 0{,}98 & 0{,}17 & 0{,}47 & 0{,}78 \\
\,+\,agent Gemini                       & 0{,}96 & 0{,}21 & 0{,}46 & 0{,}77 \\
\,+\,U-Net pancr\'eas (strict)          & 0{,}73 & 0{,}50 & 0{,}34 & 0{,}71 \\
\,+\,U-Net pancr\'eas (rel\^ach\'e)     & 0{,}92 & 0{,}30 & 0{,}45 & 0{,}77 \\
Cascade panc.$\rightarrow$tum.\ (V2)    & 0{,}94 & 0{,}40 & 0{,}50 & 0{,}80 \\
\textbf{Recall ResNet ens.}             & \textbf{0{,}88} & \textbf{0{,}74} & \textbf{0{,}54} & \textbf{0{,}81} \\
\bottomrule
\end{tabular}
\caption{Performances comparatives sur le \emph{split} test MSD
(100~images, 50 positives / 50 n\'egatives). DICE calcul\'e sur les images
positives uniquement.}
\label{tab:abl}
\end{table}

Cette observation oriente la suite de notre travail. Le segmenteur n'est
pas en cause, comme le montre la borne sup\'erieure obtenue en oracle (DICE
0{,}88 lorsqu'on fournit \`a MedSAM la bo\^ite \emph{ground-truth} de la
tumeur). Toute am\'elioration ult\'erieure doit donc s'attaquer \`a la
g\'en\'eration de bo\^ites par \emph{Grounding~DINO} ou ajouter un m\'ecanisme
de rejet en aval.

\subsection{Tentative de raffinement par agent multimodal}\label{sec:agent}

La premi\`ere id\'ee que nous explorons est de demander \`a un mod\`ele
multimodal externe de v\'erifier ou de corriger les bo\^ites douteuses. Le
notebook \texttt{agentic\_segmentation\_msd.ipynb} branche \emph{Gemini}~2.0
\cite{team2024gemini} en aval de \emph{DINO}. \`A chaque \'etape l'agent
re\c{c}oit le \emph{prompt} courant, la bo\^ite top-1, le masque \emph{MedSAM}
et un \emph{crop} de la r\'egion, et peut renvoyer une action parmi
\texttt{shrink}, \texttt{expand}, \texttt{shift}, \texttt{accept} et
\texttt{reject}.

Trois observations conduisent \`a abandonner cette piste. D'abord, l'agent
accepte la majorit\'e des hallucinations parce que rien ne distingue, dans
le \emph{crop} pr\'esent\'e, un faux positif d'un vrai. Ensuite, son
comportement reste non d\'eterministe m\^eme \`a temp\'erature nulle, ce qui
emp\^eche toute calibration fiable. Enfin, le co\^ut moyen d'environ 2~s par
appel rend l'\'evaluation prohibitive sur 100~images. Le tableau~\ref{tab:abl}
montre que ce m\'ecanisme n'apporte qu'un gain marginal de sp\'ecificit\'e
(0{,}17 $\rightarrow$ 0{,}21). Nous concluons que la v\'erification doit
provenir d'un mod\`ele plus sp\'ecialis\'e que d'un LLM g\'en\'eraliste, et
nous orientons l'effort vers des informations anatomiques explicites.

\subsection{Porte anatomique fond\'ee sur un U-Net du pancr\'eas}

Puisque les hallucinations de \emph{DINO} sont distribu\'ees sur l'ensemble
de l'image, une mani\`ere naturelle de les filtrer consiste \`a contraindre
les bo\^ites \`a tomber sur le pancr\'eas. Nous entra\^inons un \emph{U-Net}
sur le label \texttt{1} du MSD, conform\'ement \`a la section~3.2, puis nous
rejetons toute bo\^ite dont l'\emph{overlap} avec le masque pancr\'eas
dilat\'e est inf\'erieur \`a un seuil $\tau_{\text{ov}}$. Un seuil $\tau$ sur
le score \emph{DINO} compl\`ete ce filtrage spatial.

Une premi\`ere version stricte ($\tau_{\text{ov}}=0{,}10$, dilatation 15~px,
masque clip\'e au ROI) am\'eliore la sp\'ecificit\'e \`a 0{,}50 mais sacrifie
25 points de sensibilit\'e parce que le pancr\'eas pr\'edit ne couvre pas
toujours la tumeur. Une seconde version rel\^ach\'ee ($\tau_{\text{ov}}=0{,}05$,
dilatation 35~px, perte de Tversky pour le U-Net) restaure 0{,}92 de
sensibilit\'e mais ne gagne que peu en sp\'ecificit\'e (0{,}30) parce que le
ROI \'elargi laisse passer les hallucinations intra-pancr\'eatiques. La porte
anatomique appara\^it donc utile mais incompl\`ete~: elle ne sait pas
distinguer une vraie tumeur d'une fausse alarme situ\'ee sur l'organe.

\subsection{Cascade pancr\'eas$\rightarrow$tumeur exploitant le d\'etecteur}

L'observation pr\'ec\'edente sugg\`ere d'\'eviter le co\^ut d'un mod\`ele
suppl\'ementaire en exploitant directement la double sortie de notre
\emph{DINO}~v3. Comme la configuration entra\^ine simultan\'ement les classes
\texttt{pancreas} et \texttt{tumor}, nous gardons la bo\^ite \texttt{pancreas}
de plus haut score, la dilatons de 20~px, et n'acceptons une bo\^ite
\texttt{tumor} que si elle recouvre ce voisinage \`a au moins 10\,\%. Un
\emph{post-processing} \'el\'ementaire (composante connexe maximale, filtre
de taille) supprime les masques aberrants. Cette logique est document\'ee
dans \texttt{experiments/baseline\_v2/}.

Le gain est significatif (sp\'ecificit\'e 0{,}40 et DICE 0{,}50) sans entra\^iner
de mod\`ele suppl\'ementaire, ce qui simplifie la cha\^ine. La sp\'ecificit\'e
plafonne cependant en de\c{c}\`a de 0{,}50, signe qu'une part des
hallucinations partage la m\^eme localisation que les vraies tumeurs et ne
peut donc pas \^etre rejet\'ee par un simple crit\`ere spatial.

\subsection{Classifieur ResNet-18 d'ensemble}

Les sections pr\'ec\'edentes montrent que les filtres spatiaux et
s\'emantiques d\'eplacent le compromis sensibilit\'e/sp\'ecificit\'e sur la
m\^eme courbe sans en changer la forme globale. Pour modifier qualitativement
ce compromis, il faut introduire un nouveau signal d'apprentissage. Nous
entra\^inons donc un classifieur \emph{patch-level} qui juge chaque candidat
\`a partir de sa texture locale, signal qui n'est pr\'esent ni dans
\emph{DINO} ni dans le \emph{U-Net} pancr\'eas.

Le pipeline final, document\'e dans \texttt{MSD\_RECALL\_STRATEGY.md},
op\`ere ainsi. \emph{DINO} g\'en\`ere jusqu'\`a cinq candidats \`a un seuil
abaiss\'e de 0{,}05, plus permissif que pr\'ec\'edemment afin de ne pas
sacrifier la sensibilit\'e en amont. Le filtre pancr\'eas reste actif mais
relativement large (marge 35~px). Chaque candidat survivant est crop\'e en
224$\times$224 puis pass\'e \`a cinq ResNet-18 entra\^in\'es sur le m\^eme
jeu de \emph{patches} mais avec des graines diff\'erentes. Les probabilit\'es
sont agr\'eg\'ees par moyenne p\'enalis\'ee (\'equation~2). Au moment de la
d\'ecision, on conserve les candidats dont le score d'ensemble d\'epasse un
seuil $\tau_{\text{ResNet}}$ calibr\'e sur la validation par balayage
$F_{\beta=2}$ avec contrainte de pr\'ecision $P\geq0{,}80$. Jusqu'\`a deux
candidats sont alors envoy\'es \`a \emph{MedSAM} et leurs masques sont
fusionn\'es par union.

La qualit\'e de cette d\'ecision repose enti\`erement sur le jeu d'entra\^inement
du classifieur. Si l'on se contentait de positifs et de n\'egatifs choisis au
hasard, le ResNet apprendrait des indices triviaux comme la pr\'esence de
muqueuse intestinale ou la teinte g\'en\'erale, et n'attaquerait pas le
probl\`eme r\'eel des hallucinations de \emph{DINO}. Nous \'echantillonnons
donc les n\'egatifs durs en abaissant le seuil de \emph{DINO} \`a 0{,}01 et
$\textit{top-k}=12$ pendant l'extraction (\texttt{extract\_hard\_negatives.py}),
suivant l'esprit d'\emph{Online Hard Example Mining}~\cite{shrivastava2016ohem}.
Les positifs \emph{ground-truth} sont l\'eg\`erement \emph{jitter\'es}
(translation $\pm$12\,\%, \emph{scale} $\pm$18\,\%) pour robustifier le
mod\`ele, et le ratio final est born\'e entre 2:1 et 4:1 n\'egatifs/positifs.

Le tableau~\ref{tab:abl} reporte le gain. La sp\'ecificit\'e atteint 0{,}74
contre 0{,}17 pour le \emph{baseline}, soit une am\'elioration de 57 points,
au prix d'une perte de 10 points de sensibilit\'e. Le F1 monte \`a 0{,}81 et
le DICE moyen sur images positives \`a 0{,}54.

\subsection{Calibration du seuil et analyse des faux n\'egatifs}

Les seuils \emph{per-fold} obtenus pendant l'entra\^inement (tableau~\ref{tab:thresh})
montrent une variabilit\'e marqu\'ee, signe que les cinq r\'eseaux apprennent
des d\'ecisions sensiblement diff\'erentes. Cette variabilit\'e justifie
\emph{a posteriori} le choix d'un ensemble plut\^ot que d'un mod\`ele unique.
Le seuil global retenu pour le test (0{,}37) provient de l'agr\'egation par
moyenne p\'enalis\'ee, qui d\'eplace l\'eg\`erement la d\'ecision vers la
prudence.

\begin{table}[t]
\centering
\small
\begin{tabular}{lr}
\toprule
\emph{Fold} ResNet & Seuil $\tau_{\beta=2}$ \\
\midrule
1 & 0{,}539 \\
2 & 0{,}343 \\
3 & 0{,}564 \\
4 & 0{,}841 \\
5 & 0{,}641 \\
\midrule
M\'ediane & 0{,}564 \\
Seuil ensemble retenu & 0{,}370 \\
\bottomrule
\end{tabular}
\caption{Seuils ResNet par \emph{fold} sur la validation. La large dispersion
motive l'agr\'egation par moyenne p\'enalis\'ee.}
\label{tab:thresh}
\end{table}

Le script \texttt{test\_gd\_msd\_final\_recall.py} \'etiquette chaque faux
n\'egatif selon sa cause directe. La majorit\'e tombe dans la cat\'egorie
\texttt{dino\_or\_spatial\_filters}, c'est-\`a-dire que \emph{DINO} ne propose
aucun candidat sur la coupe concern\'ee. Le classifieur n'est donc pas
limitant en sensibilit\'e, et la borne sup\'erieure que pourrait atteindre
notre cha\^ine d\'epend principalement du \emph{recall} du d\'etecteur
amont.

\section{Analyse Critique}

Le \emph{fine-tuning} de \emph{Grounding~DINO} avec \emph{focal loss} et
trois cha\^ines \emph{distractor} se r\'ev\`ele indispensable~; sans cet
ajustement, le mod\`ele transf\'er\'e en \emph{zero-shot} ne d\'epasse pas
mAP$<$0{,}05 et la cha\^ine compl\`ete devient inutilisable. La cha\^ine
\emph{detector$+$MedSAM} reste un \emph{baseline} solide en sensibilit\'e
d\`es lors que le d\'etecteur est sp\'ecialis\'e, parce que \emph{MedSAM}
tol\`ere une bo\^ite m\^eme grossi\`ere et atteint 0{,}88 de DICE en oracle.

L'utilisation de \emph{prompts} n\'egatifs tels que ``\texttt{no~tumor}''
\'echoue parce que l'encodeur BERT~\cite{devlin2019bert} ne traite pas la
n\'egation s\'emantiquement. L'introduction de \emph{distractors}
positifs (``\texttt{normal pancreas tissue}'', ``\texttt{bowel gas}'') a un
effet plus tangible mais limit\'e \`a une dizaine de points de pr\'ecision.

L'\'echec de la boucle agentique m\'erite quelques mots. Nous esp\'erions
qu'un mod\`ele multimodal \emph{frontier} pourrait jouer le r\^ole d'un
second avis radiologique, mais trois facteurs combinent leurs effets contre
cette esp\'erance. Le \emph{crop} pr\'esent\'e a une r\'esolution insuffisante
pour discriminer une tumeur d'un parenchyme sain, le mod\`ele tend \`a
accepter par d\'efaut, et son non-d\'eterminisme rend toute calibration peu
fiable.

L'analyse comparative du tableau~\ref{tab:abl} d\'egage un sch\'ema clair.
Les portes purement spatiales ou par seuil d\'eplacent le point de
fonctionnement le long d'une m\^eme courbe sensibilit\'e/sp\'ecificit\'e
sans changer sa forme. Seul le ResNet d'ensemble parvient \`a d\'ecaler la
courbe globalement, parce qu'il apporte un signal d'apprentissage suppl\'ementaire
que ne capture aucun autre composant. Sa d\'ecomposition en \emph{region
proposal} amont et classification \emph{patch-level} aval rappelle
\emph{Mask R-CNN}~\cite{he2017maskrcnn}, et confirme empiriquement que la
texture locale d'un \emph{patch} est un crit\`ere compl\'ementaire \`a un
score d'alignement texte-image.

Notre protocole d'\'evaluation pr\'esente n\'eanmoins deux limites
m\'ethodologiques importantes. D'une part, toutes les m\'etriques sont
calcul\'ees sur des coupes 2D ind\'ependantes, alors que les tumeurs
pancr\'eatiques apparaissent typiquement sur plusieurs coupes adjacentes,
contrairement aux hallucinations. Une simple agr\'egation \emph{multi-slice}
am\'eliorerait probablement la sp\'ecificit\'e sans co\^ut sensibilit\'e.
D'autre part, le \emph{split} 50/50 du test a \'et\'e construit pour faciliter
la comparaison entre m\'ethodes, mais ne refl\`ete pas la pr\'evalence
clinique r\'eelle des coupes pathologiques. Une \'evaluation sur une
distribution naturelle r\'e\'evaluerait la valeur pratique de chaque
composant.

Enfin, l'ensemble des \emph{checkpoints} (DINO~v3, U-Net pancr\'eas, cinq
ResNet-18) tient en moins d'un gigaoctet, et l'inf\'erence \emph{end-to-end}
sur une coupe prend environ 300~ms sur GPU~T4. Le co\^ut total d'entra\^inement
reste sous 90~minutes de calcul GPU, soit un ordre de grandeur inf\'erieur
\`a celui de \emph{nnU-Net} 3D pour un probl\`eme comparable, ce qui rend
notre approche int\'eressante pour des d\'eploiements rapides.

\section{Conclusion}

Notre projet documente une transition. Sur \emph{Kvasir-SEG}, la cha\^ine
\emph{Grounding~DINO}+\emph{MedSAM} fonctionne pratiquement sans intervention
humaine et atteint un DICE oracle de 0{,}95. Sur \emph{MSD~Pancreas}, la
m\^eme cha\^ine s'effondre en sp\'ecificit\'e, et seule l'introduction d'un
classifieur appris en aval restaure un compromis exploitable. Nos
exp\'eriences montrent que les filtres spatiaux et les agents multimodaux
ne d\'eplacent que la position le long d'une courbe fixe, tandis qu'un
signal d'apprentissage \emph{patch-level}, entra\^in\'e sur les n\'egatifs
durs g\'en\'er\'es par le d\'etecteur lui-m\^eme, change qualitativement le
compromis. Le pipeline final atteint un F1 de 0{,}81 et un DICE moyen de
0{,}54 sur les coupes positives. L'int\'egralit\'e du code, des donn\'ees et
des \emph{checkpoints} est publi\'ee \`a~\url{https://github.com/moradBMH/INF8225_Projet}.

\section*{Remerciements}

Nous remercions les charg\'es de laboratoire INF8225 pour leurs retours
m\'ethodologiques, ainsi que les auteurs de \emph{MedSAM},
\emph{Grounding~DINO} et de la base \emph{Medical Segmentation Decathlon}
pour la mise \`a disposition de leurs artefacts.

\bibliographystyle{named}
\bibliography{ijcai19}

\end{document}
